% SPDX-License-Identifier: Apache-2.0
%
% The TxHDL cheat sheet: two pages, landscape, meant to be printed on
% the two sides of one sheet. The front walks one small unit from its
% source to its netlist and its waveform; the back is the vocabulary.
% Each side stands on its own. Built by //docs:cheatsheet, and as one
% PNG per side by //docs:cheatsheet_png.
\documentclass[10pt]{article}
\usepackage[paperwidth=11in,paperheight=8.5in,margin=0.4in,top=0.35in,bottom=0.35in]{geometry}
\usepackage{multicol}
\usepackage{array}
\newcommand{\listingsize}{\fontsize{6.2}{7.2}\selectfont}
% SPDX-License-Identifier: Apache-2.0
%
% The house preamble, shared by every document in this directory. Loaded
% after \documentclass, before \title. Keeping it in one file is what
% stops two articles drifting apart in font, listing style or figure
% style.
% Computer Modern. IEEEtran selects Times (ptm) by default, so the face has
% to be chosen explicitly.
%
% Latin Modern is the Computer Modern variety used here, and the choice is
% forced rather than aesthetic. Under T1 encoding, plain `cmr` has no Type 1
% outlines unless cm-super is installed, and it is not in the pinned
% distribution, so pdflatex falls back to EC bitmaps and every glyph in the
% PDF becomes a Type 3 font. That was measured: the first build produced a
% document with 10 Type 3 fonts and no Type 1. Latin Modern is Computer
% Modern redrawn with a full T1 repertoire and real .pfb outlines, so the
% same design comes out scalable.
\usepackage[T1]{fontenc}
\usepackage[utf8]{inputenc}
\usepackage{lmodern}
% microtype is not in the pinned TeX distribution. emergencystretch alone
% keeps the two-column measure from producing overfull lines.
\setlength{\emergencystretch}{3em}

\usepackage{amsmath}
\usepackage{amssymb}
\usepackage{amsthm}
\usepackage{booktabs}
\usepackage{array}
% enumitem is not in the pinned TeX distribution; the list spacing below
% does what \setlist would have done.
\usepackage{float}
\usepackage{xcolor}
\usepackage{listings}
\usepackage{tikz}
\usetikzlibrary{arrows.meta,positioning,fit,backgrounds,calc}
\usepackage[hidelinks,bookmarks=true]{hyperref}

% Numbered, definition-style box for a rule the reader is meant to apply.
\theoremstyle{definition}
\newtheorem{rulebox}{Rule}
\newtheorem{example}{Example}

% Unnumbered asides.
\theoremstyle{remark}
\newtheorem*{tip}{Tip}
\newtheorem*{pitfall}{Pitfall}

% Tighter lists, in place of enumitem's \setlist.
\setlength{\topsep}{2pt}
\setlength{\itemsep}{1pt}
\setlength{\parsep}{0pt}

\newcommand{\code}[1]{\texttt{#1}}
\newcommand{\term}[1]{\emph{#1}}

% Syntax highlighting. Four roles, distinguishable in colour and still
% distinguishable in greyscale, because an IEEE article gets printed: the
% keyword is the only bold role, the comment is the only italic one, and the
% two remaining roles differ in lightness as well as in hue.
\definecolor{kwblue}{RGB}{20,60,150}
\definecolor{cmtgreen}{RGB}{90,120,90}
\definecolor{strred}{RGB}{150,50,40}
\definecolor{numgray}{gray}{0.45}
\definecolor{framegray}{gray}{0.70}
\definecolor{bgshade}{gray}{0.97}
% A darker fill for figure cells. The listing background has to stay very
% light to keep code readable; a figure cell has to be clearly shaded or the
% distinction it draws is invisible in print.
\definecolor{cellshade}{gray}{0.90}

% The listing size is the document's choice, because it depends on the
% column: a two-column article sets `\listingsize` to 6pt and keeps its
% listings within 70 columns, a one-column document keeps the body size
% and includes source files at 80. Lines are never broken; a line that does not fit
% overflows the frame, which is visible, rather than wrapping, which is
% not.
\providecommand{\listingsize}{\scriptsize}

% `'` is deliberately not a string delimiter: the language uses it for loop
% labels, as Rust does, so treating it as a quote would swallow the rest of
% the line.
\lstdefinestyle{txhdl}{
  basicstyle=\ttfamily\listingsize,
  keywordstyle=\bfseries\color{kwblue},
  commentstyle=\itshape\color{cmtgreen},
  stringstyle=\color{strred},
  numberstyle=\tiny\color{numgray},
  morekeywords={mod,use,pub,const,type,struct,enum,trait,impl,fn,async,let,mut,
    match,if,else,loop,while,for,in,break,continue,return,as,await,
    module,bus,channel,transaction,protocol,pipeline,flow,seq,config,
    port,stage,pipe,instance,connect,bind,tag,domain,blackbox,foreign,
    property,role,signal,handler,true,false,
    sequence,interface,modport,implements,package,import,var,wire,func,proc,
    case,default,next,fork,branch,join,state,fsm},
  morecomment=[l]{//},
  morestring=[b]",
  showstringspaces=false,
  columns=fullflexible,
  keepspaces=true,
  breaklines=false,
  % Framed, so a listing is bounded rather than floating in the column.
  frame=single,
  rulecolor=\color{framegray},
  framesep=3pt,
  backgroundcolor=\color{bgshade},
  xleftmargin=3pt,
  xrightmargin=3pt,
  aboveskip=6pt,
  belowskip=6pt,
  % Every listing is numbered and captioned, so it can be referred to by
  % number. `caption` without `float` numbers the listing and leaves it
  % where it was written, which is what a listing in running prose needs.
  captionpos=b,
  abovecaptionskip=2pt,
  belowcaptionskip=6pt,
}

% Figures drawn as text. Framed the same way, and with the highlighting
% switched off, because a box-and-arrow diagram is not code and half its
% words turning blue reads as an accident. Setting a style adds keys rather
% than clearing the ones already set, so the three roles are neutralised by
% hand rather than by leaving them out.
\lstdefinestyle{ascii}{
  basicstyle=\ttfamily\listingsize,
  keywordstyle=\ttfamily,
  commentstyle=\ttfamily,
  stringstyle=\ttfamily,
  columns=fullflexible,
  keepspaces=true,
  frame=single,
  rulecolor=\color{framegray},
  framesep=3pt,
  backgroundcolor=\color{bgshade},
  xleftmargin=3pt,
  xrightmargin=3pt,
  aboveskip=6pt,
  belowskip=6pt,
}
\lstset{style=txhdl}

% House TikZ styles. `shorten` lives on the arrow styles rather than on each
% arrow, so every arrow in the document gets the gap, including ones added
% later. TikZ clips a path at a node's border, so without it an arrowhead
% butts against the box with no gap at all. Keep `node distance` well above
% twice the shorten, or the arrow shrinks to nothing.
\tikzset{
  box/.style={draw,rounded corners=2pt,align=center,inner sep=4pt,
              font=\footnotesize},
  boxf/.style={box,fill=bgshade},
  lbl/.style={font=\scriptsize\itshape,align=center},
  fl/.style={-{Stealth[length=4pt]},shorten >=2.5pt,shorten <=2.5pt},
  fld/.style={fl,dashed},
}


\pagestyle{empty}
\setlength{\columnsep}{16pt}
\setlength{\parindent}{0pt}
\setlength{\parskip}{2pt}
\definecolor{bar}{RGB}{20,60,150}
\definecolor{tint}{RGB}{232,238,250}

% A section head: a coloured bar.
% A heading never sits alone at the foot of a column: when less than
% five lines are left, the column breaks before it.
\newcommand{\keep}{\par\vskip 0pt plus 5\baselineskip\penalty-200\vskip 0pt plus -5\baselineskip}
\newcommand{\head}[1]{\keep\vspace{4pt}\noindent\colorbox{bar}{\parbox{\dimexpr\linewidth-2\fboxsep}{\color{white}\bfseries\small #1}}\vspace{2pt}\par}

% The banner across the top of a side.
\newcommand{\banner}[2]{\noindent\colorbox{bar}{\parbox{\dimexpr\textwidth-2\fboxsep}{\color{white}
{\LARGE\bfseries #1}\hfill{\small #2}}}\vspace{4pt}\par}

% A two-column table of code and what it means, a column wide.
\newenvironment{vocab}{\par\footnotesize\setlength{\tabcolsep}{3pt}\begin{tabular}{@{}>{\raggedright\arraybackslash}p{0.53\linewidth}>{\raggedright\arraybackslash}p{0.44\linewidth}@{}}}{\end{tabular}\par}

\begin{document}
\small

% ---------------------------------------------------------------- front
\banner{TxHDL cheat sheet}{hardware in Rust: the front walks one unit
through; the back is the vocabulary}

\begin{multicols}{3}

\head{The idea}
You already know most of TxHDL, because it is Rust. Where Rust has a
way to say something, TxHDL uses it and adds nothing: a transaction
is a \code{struct}, a unit is a \code{struct} with an \code{async fn
run}, a pipeline is an \code{async fn}, a clock is a type. So
\code{rustc} does the checking for you. Two drivers on one wire is a
move error. A signal handed to the wrong clock is a type error. A
cycle count in combinational code can't happen, because there is no
\code{.await} to write.

\head{The one rule about time}
A process \textbf{waits for an edge, then reads, then drives}.
Every \code{.await} is a cycle boundary. After the wait, a register
reads as the edge left it, and anything you drive lands at the end
of the step, all at once. That is why you can read and drive the
same register in one step and still mean hardware.

The waits you'll use: \code{C::rising()}, \code{C::falling()}, and
\code{until(C::rising, || cond)} for ``the first edge where
\code{cond} holds''. Two events on one edge are
\code{parallel!(a, b).await}. Time counts in ticks, two per cycle of
the default clock, so a falling edge has a tick of its own.

\head{Build it, run it, look at it}
All you install is \code{bazelisk}. The toolchains, the simulators
and the \LaTeX{} for this page come through the build.
\begin{lstlisting}[style=ascii,frame=none,backgroundcolor=\color{tint},xleftmargin=2pt]
bazel build //...                 # everything
bazel test  //...                 # and check it
bazel run //lib/examples:ex_cheat # the netlist
TXHDL_FST=c.fst bazel run \
  //lib/examples:ex_cheat         # and a waveform
bazel build //docs/...            # every document
\end{lstlisting}

\head{Where to read more}
\code{//docs:tutorial} takes you through in nine steps, one example
each. \code{//docs:examples} has one example per idea, with what each
printed. \code{//docs:zero} starts from an empty directory.
\code{//docs:cover} lists every document there is.

\head{Adding your own}
Copy \code{lib/examples/ex\_cheat.rs} to \code{ex\_\emph{name}.rs},
add it to \code{lib/examples/BUILD.bazel} and its \code{sources},
give it a \code{waveform()} entry, and write its section in the
examples document. An idea lives in three places: the library, an
example, and a document. The example is the test.


\columnbreak

\head{A whole unit: a counter}
This counts while \code{enable} is high and raises \code{tick} on
seven. Read it top to bottom: the struct is the state, the
\code{impl} is the behaviour, and the ports are in the signature.
\lstinputlisting[rangebeginprefix=//\ begin\{,rangebeginsuffix=\},rangeendprefix=//\ end\{,rangeendsuffix=\},includerangemarker=false,linerange=unit-unit,frame=none,backgroundcolor=\color{tint},xleftmargin=2pt]{lib/examples/ex_cheat.rs}
\code{\#[lower]} is what makes the same code a netlist too. Leave it
off and the unit still simulates.

\head{What it turns into}
Here is what \code{Counter::verilog("counter")} printed when the
build ran it. The \code{rst} is the reset every module gets for free.
The build also writes the VHDL, and simulates both against the Rust
run on the right, tick by tick.
\lstinputlisting[style=ascii,frame=none,backgroundcolor=\color{tint},xleftmargin=2pt]{out_cheat.txt}

\columnbreak

\head{And what it did}
The build drew this from the waveform the run wrote, so it can't
disagree with the code. \code{enable} drops at cycle 10, \code{n}
counts while it's high, and \code{tick} marks each seven.
\begin{center}
\input{cheat_timing_color.tex}
\end{center}

\head{The shape every unit has}
\begin{itemize}\setlength{\itemsep}{0pt}\setlength{\leftmargin}{8pt}
\item One \code{loop}, and its first line is the wait.
\item Then reads: inputs with \code{get}, registers just by name.
\item Then \code{let}s, each a named wire in the netlist.
\item Then drives, each written once: \code{with!} for registers,
  \code{set} for outputs, \code{send} for channels.
\end{itemize}
Nothing is assigned twice and nothing branches around a drive. If
you keep to that shape, the lowering can read your unit.

\head{How the netlist is checked}
For every lowered example the build keeps the Rust run's trace,
writes VHDL and Verilog, and replays the inputs into both, asserting
every register and output at every tick: VHDL under nvc and Verilog
under Verilator. Add \code{lowered = (entity, unit)} to a
\code{waveform()} in \code{docs/BUILD.bazel} and your unit gets the
same two tests.


\head{When something goes wrong}
The macros refuse what they can't lower, at compile time, with a
message that names the construct. If the Rust run and the netlist
ever disagree, the co-simulation test fails at the first tick where
they differ. Look at the waveform first; it usually tells you why.

\vfill
\InputIfFileExists{buildstamp.tex}{}{}
\providecommand{\buildstamp}{Draft build (outside the Bazel flow).}
{\footnotesize\buildstamp\ Written with a large language model,
Claude, as an assistant, as every commit says. Turn over for the
vocabulary.}

\end{multicols}

% ---------------------------------------------------------------- back
\newpage
\banner{TxHDL, the vocabulary}{everything you'll reach for, and how it
lowers}
\footnotesize
% The back holds the whole vocabulary, so its tables are a size smaller.
\renewenvironment{vocab}{\par\scriptsize\setlength{\tabcolsep}{3pt}\begin{tabular}{@{}>{\raggedright\arraybackslash}p{0.55\linewidth}>{\raggedright\arraybackslash}p{0.42\linewidth}@{}}}{\end{tabular}\par}

\begin{multicols}{3}

\head{Values and state}
\begin{vocab}
\code{Bit}, \code{U<N>}, \code{I<N>} & bits and words; Rust's operators work on them \\
\code{\#[derive(Value)]} & your own enum or struct as a value, shown by name in a trace \\
\code{\#[derive(Transaction)]} & a struct that can travel on a channel \\
\code{Reg<T, C>} & a register on clock \code{C} (default if left out); read it by name \\
\code{Wire<T>} & a \code{let} kept as a field, so the waveform can see it \\
\code{Mem<T, N>} & a memory; \code{read(a)}, and write with \code{at(a)} in \code{with!} \\
\code{slice}, \code{concat}, \code{sext}, \code{zext} & take words apart, put them together \\
\code{mul::<M>()} & a product wider than its operands \\
\code{U<256, 2>} & wider than 128 bits: the second number is how many 128-bit limbs \\
\end{vocab}

\head{Driving}
\begin{vocab}
\code{with!(self <= \{ c ? n: v \})} & drive registers, each under its condition \\
\code{when!(c => self \{ .. \})} & the same, condition first \\
\code{case!(v => \{ p => \{..\} \})} & many arms; the first match wins \\
\code{if c \{ .. \}} & Rust's own \code{if}, around drives \\
\code{out.set(v)} & drive an output wire this cycle \\
\end{vocab}

\head{Choosing a value}
\begin{vocab}
\code{mux(c, a, b)} & two ways \\
\code{select!(v => \{ p => e, .. \})} & many ways, by pattern \\
\code{if c \{ a \} else \{ b \}} & Rust's form, lowered as a \code{mux} \\
\code{match v \{ .. \}} & lowered as a \code{select!} \\
\code{0x30..=0x39}, \code{(lo..=hi).contains(\&x)} & ranges, as patterns or tests \\
\end{vocab}
Hardware has no branch: both sides exist and the condition picks.
The names are there so the code reads the way the circuit works.

\head{Wires and channels}
\begin{vocab}
\code{signal()} $\to$ \code{(Out, In)} & a wire's two ends; \code{Out} can't be cloned, so one driver \\
\code{chan()} $\to$ \code{(Tx, Rx)} & a channel with a buffer of two \\
\code{tx.send(v)}, \code{tx.ready()} & offer a word; is there room \\
\code{rx.peek()}, \code{rx.recv()} & look at the head; take it \\
\code{rx.recv\_if(c)}, \code{rx.take()} & take under a condition; take whatever is there \\
\code{pad()} & a pin driven from both sides \\
\end{vocab}
In the netlist a channel is its data, a \code{valid} and a
\code{ready}. A wait on \code{until} becomes the guard of everything
after it.

\head{Ports with names}
One port per side is fine for a counter. For more, give each side a
struct and name the fields, so two ports of the same type can't be
swapped by accident:
\begin{vocab}
\code{run(\&mut self, i: DiffIn, o: DiffOut)} & each side any struct of ports \\
\code{DiffIn \{ minuend, subtrahend \}: DiffIn} & take it apart right in the signature \\
\code{\#[derive(Ports)]} & all a struct from another crate needs \\
a field that is itself a struct of ports & its ports under the field's name, \code{left\_aw\_valid} \\
\code{side: [Rx<T>; N]} & an array of ports, \code{side\_0} onward \\
\code{interface! \{ Ops \{..\} role Source \{..\} \}} & declare a link once, get a struct per role \\
\end{vocab}
The netlist flattens each struct into its fields, under their names.
Every AXI-Lite peripheral takes its link as one struct,
\code{LitePort}.

\head{Units of units}
A parent keeps its children as fields, hands each its ports in
\code{run}, and runs them with \code{join2} (nest it for more). The
lowering writes one module with an instance per child.
\begin{vocab}
\code{let (h, p) = link::<LitePort<..>>()} & a whole link's channels in one line \\
\code{child.run((x, tie(v)), y)} & an input held at a constant \\
\code{child.run((bus, ..), ..)} & a side passed on whole, from any crate \\
\end{vocab}

\head{Clocks}
\begin{vocab}
\code{struct Slow; impl Clock for Slow} & a clock: a name and a period in ticks \\
\code{Reg<T, Slow>}, \code{In<T, Slow>} & state and ports on that clock \\
\code{Slow::rising().await} & a process waits on one clock \\
\code{Crossing<T, A, B>} & the only way a signal changes clock \\
\code{ChanCdc} & a channel between clocks, in \code{//lib/parts} \\
\end{vocab}
A design with two clocks is a parent with a child on each; the parent
names neither. \code{DefaultClock} you never have to name at all.

\head{Reset}
You don't write one. Every module gets \code{rst}, and while it's
high every register returns to its first value and every channel
empties. If a unit must \emph{do} something on reset, take
\code{rst: In<Bit>} and read it; it joins the same net.

\head{Other people's hardware}
A hand-written Verilog or VHDL module can be a child: give it a Rust
model for the simulation and an \code{impl Lower} that returns
\code{netlist::foreign(..)}. The netlist instantiates it by its own
name. Going the other way, the build rules \code{verilog\_unit()}
and \code{vhdl\_unit()} turn a netlist back into a unit you can run
next to the Rust it came from.

\head{What the lowering reads}
A loop of one wait, then reads, \code{let}s and drives: one clocked
block. Several waits, or one under \code{if}, make a state machine,
a state per wait, and an output a state sets holds until another
sets it. \code{for i in 0..N} over an array of ports is unrolled for
each \code{N}. Inside,
Rust's operators: \code{+ - * \% \& | \^{} ! <{}< >{}>}, unary \code{-},
the compares, and \code{as} to an unsigned type. Calls to a
\code{fn} marked \code{\#[lower]} in the same file are inlined. What
you can't have is what hardware can't: a call the lowering doesn't
know, or a loop whose count varies.

\head{Names}
Call things what they are. A port named \code{next} or \code{begin},
or a register named \code{signal}, is reserved in VHDL or Verilog;
the netlist writes it as \code{next\_rw}, the same in both, and the
testbench follows. Want a nicer name in the netlist?
\code{\#[rename("nxt")]} on the field.

\head{Register maps and checks}
\begin{vocab}
\code{regmap! \{ m (m\_read, m\_we), 2: [..] \}} & a peripheral's map said once: the decode, the C header, the datasheet table \\
\code{(1, ctrl, rw, "..", [(run, 0, 1, ..)])} & a register and its bit fields \\
\code{check!(c, "why")} & must hold where it is reached; an \code{assert} in the netlist \\
\code{assume!(c, "why")}, \code{cover!(c, "why")} & taken for granted; can happen \\
\end{vocab}

\head{Pipelines and builds}
\begin{vocab}
\code{\#[pipeline(..)] async fn} & a pipeline; stages fall where you \code{.await} \\
\code{pipeline::drive(f, rx, tx)} & one call of \code{f} per cycle \\
\code{config! \{ Fpga: C for Top<Fpga> \{..\} \}} & a build as one item \\
\end{vocab}

\head{Watching and testing}
\begin{vocab}
\code{\#[derive(Trace)]} & fields named in the waveform \\
\code{Wave::from\_env()} & a writer when \code{TXHDL\_FST} or \code{TXHDL\_VCD} is set \\
\code{w.add("name", \&x)}, \code{start()}, \code{stop()} & choose what to watch \\
\code{Running::new(u.run(..))} & start the design \\
\code{sim.cycle()} & one cycle; look at things between calls \\
\end{vocab}
A unit runs in a plain \code{\#[test]} too, each test with its own
time. In a test, compare registers after a cycle rather than output
wires.

\end{multicols}
\end{document}
